\documentclass[../../main.tex]{subfiles}% 注意这里的文件路径不能用 ./main.tex ,否则用latexmk编译子文件会报错
\graphicspath{{\subfix{./image/}}{\subfix{../../image/}}} % 指定图片引用路径,后续可以直接使用图片文件名
% 如果图片存放在上级目录,则使用 ../../image/ 作为备用路径

% 例如:
% \begin{figure}[H]
% \centering
% \includegraphics[width=0.3\textwidth]{图.png}
% \caption{}
% \label{figure:图}
% \end{figure}
% 注意:上述\label{}一定要放在\caption{}之后,否则引用图片序号会只会显示??.

\begin{document}

\section{收敛的概念}\label{section:ch05-Radon弱星}

\subsection{强收敛}

\cref{proposition:ch02-强收敛性质}启发我们给出以下定义，特别地，当 $X$ 为赋 Borel $\sigma$-代数的局部紧 Hausdorff 空间时，该定义适用于 $\mathscr{R}_f(X)$。

\begin{definition}\label{definition:ch05-强收敛}
设 $(X,\mathcal{A})$ 为可测空间。序列 $\{\mu_k\}\subseteq\mathcal{M}_f(X)$ \textbf{强收敛}于 $\mu\in\mathcal{M}_f(X)$（记作 $\mu_k\to\mu$），若
\begin{align*}
\lim_{k\to\infty}\|\mu_k-\mu\|
=\lim_{k\to\infty}|\mu_k-\mu|(X)=0.
\end{align*}
$\mathcal{M}_f(X)$ 上的强拓扑为与范数 $\|\cdot\|$ 相伴的度量拓扑。为清楚起见，我们将上述范数在子空间 $\mathscr{R}_f(X)\subseteq\mathcal{M}_f(X)$ 上的限制记作 $\|\cdot\|_{\mathscr{R}_f(X)}$。
\end{definition}

\subsection{弱$^*$收敛}

\cref{theorem:ch02-Rf对偶}启发我们给出以下定义。

\begin{definition}\label{definition:ch05-弱星收敛}
设 $X$ 为局部紧 Hausdorff 空间。序列 $\{\mu_k\}\subseteq\mathscr{R}_f(X)$ \textbf{弱$^*$收敛}于 $\mu\in\mathscr{R}_f(X)$（记作 $\mu_k\stackrel{*}{\rightharpoonup}\mu$），若
\begin{align}
\lim_{k\to\infty}\int_X g\,\mathrm{d}\mu_k=\int_X g\,\mathrm{d}\mu
\quad\text{对任意 }g\in C_0(X).\label{align:ch05-5.1}
\end{align}
$\mathscr{R}_f(X)$ 上的弱$^*$拓扑，记作 $\mathcal{T}_{w^*}$，是使所有映射 $\mu\mapsto\int_X g\,\mathrm{d}\mu$（$g\in C_0(X)$）连续的最弱拓扑。
\end{definition}

\begin{remark}\label{remark:ch05-弱星积分}
由于 $C_0(X)\subseteq C_b(X)\subseteq L^1(X,\mathscr{B}(X),|\mu|)$ 且对任意 $k\in\mathbb{N}$ 有 $C_0(X)\subseteq C_b(X)\subseteq L^1(X,\mathscr{B}(X),|\mu_k|)$，故\eqref{align:ch05-5.1}与下文\eqref{align:ch05-5.3}中的积分是良定义的（见\cref{definition:ch02-Lebesgue}）。
\end{remark}

$\mathcal{T}_{w^*}$ 的一组基为 $\mathscr{R}_f(X)$ 的如下子集族：
\begin{align}
B_{g_1,\cdots ,g_n;\epsilon}(\mu_0)
\triangleq \left\{\mu\in\mathscr{R}_f(X)\ \middle|\
\left|\int_X g_j\,\mathrm{d}\mu-\int_X g_j\,\mathrm{d}\mu_0\right|<\epsilon
\ \forall j=1,\cdots ,n\right\},\label{align:ch05-5.2}
\end{align}
其中 $g_1,\cdots ,g_n\in C_0(X)$，$\mu_0\in\mathscr{R}_f(X)$，$n\in\mathbb{N}$，$\epsilon>0$。

在对 $X$ 的适当假设下，限制在 $\mathscr{R}_f(X)$ 有界子集上的弱$^*$拓扑 $\mathcal{T}_{w^*}$ 是可度量化的：

\begin{proposition}\label{proposition:ch05-弱星可度量化}
\label{theorem:ch05-弱星可测}
设 $X$ 为具可数基的局部紧 Hausdorff 空间，$M\subseteq\mathscr{R}_f(X)$ 有界。则相对弱$^*$拓扑 $\mathcal{T}_{w^*}\cap M$ 可度量化。
\end{proposition}

\begin{proof}
空间 $C_0(X)$ 赋以 $\|\cdot\|_\infty$ 范数是可分的（见附录 A,\cref{subsection:appA-可分})。于是由\cref{theorem:ch02-Rf对偶}与\cref{proposition:appA-弱星可度量化}即得结论。
\end{proof}

\subsection{窄收敛}

在\cref{definition:ch05-弱星收敛}中选取不同的试验函数类，便得到 $\mathscr{R}_f(X)$ 上另一种收敛概念。

\begin{definition}\label{definition:ch05-窄收敛}
设 $X$ 为局部紧 Hausdorff 空间。序列 $\{\mu_k\}\subseteq\mathscr{R}_f(X)$ \textbf{窄收敛}于 $\mu\in\mathscr{R}_f(X)$（记作 $\mu_k\stackrel{n}{\rightharpoonup}\mu$），若
\begin{align}
\lim_{k\to\infty}\int_X g\,\mathrm{d}\mu_k=\int_X g\,\mathrm{d}\mu
\quad\text{对任意 }g\in C_b(X).\label{align:ch05-5.3}
\end{align}
$\mathscr{R}_f(X)$ 上的窄拓扑，记作 $\mathcal{T}_n$，是使所有映射 $\mu\mapsto\int_X g\,\mathrm{d}\mu$（$g\in C_b(X)$）连续的最弱拓扑。
\end{definition}

窄拓扑被某些作者称为“弱拓扑”（例如见文献~\cite{ref45}）。$\mathcal{T}_n$ 的一组基为\eqref{align:ch05-5.2}所定义的子集族 $B_{g_1,\cdots ,g_n;\epsilon}(\mu_0)\subseteq\mathscr{R}_f(X)$，只是其中 $g_1,\cdots ,g_n\in C_b(X)$。

\begin{remark}\label{remark:ch05-模糊收敛}
序列 $\{\mu_k\}\subseteq\mathscr{R}_f(X)$ \textbf{模糊收敛}于 $\mu\in\mathscr{R}_f(X)$，定义为
\begin{align*}
\lim_{k\to\infty}\int_X g\,\mathrm{d}\mu_k=\int_X g\,\mathrm{d}\mu
\quad\text{对任意 }g\in C_c(X)
\end{align*}
（例如见文献~\cite{ref42}）。显然，在 $\mathscr{R}_f(X)$ 中，强收敛 $\Rightarrow$ 窄收敛 $\Rightarrow$ 弱$^*$收敛 $\Rightarrow$ 模糊收敛。
\end{remark}

\begin{remark}\label{remark:ch05-范数下半连续}
\label{remark:ch05-5.1.3}
由范数的下半连续性可知，若 $\mu_k\stackrel{*}{\rightharpoonup}\mu$（特别地，若 $\mu_k\stackrel{n}{\rightharpoonup}\mu$），则
\begin{align}
|\mu|(X)=\|\mu\|_{\mathscr{R}_f(X)}
\leqslant\liminf_{k\to\infty}\|\mu_k\|_{\mathscr{R}_f(X)}
=\liminf_{k\to\infty}|\mu_k|(X).\label{align:ch05-5.4}
\end{align}
特别地，若 $\{\mu_k\}\subseteq\mathscr{R}^+_f(X)$，则 $\mu(X)\leqslant\liminf_{k\to\infty}\mu_k(X)$。具体地，若 $\mu_k\stackrel{n}{\rightharpoonup}\mu$，则在\eqref{align:ch05-5.3}中取 $g=1$ 得 $\lim_{k\to\infty}\mu_k(X)=\mu(X)$，故“没有质量损失”。若 $\mu_k\stackrel{*}{\rightharpoonup}\mu$，则未必如此：取 $X=\mathbb{R}$，$\mu_k=\delta_k$（$k\in\mathbb{N}$），则 $\mu_k\stackrel{*}{\rightharpoonup}0$ 且对所有 $k$ 有 $\mu_k(\mathbb{R})=1$。故在此情形不等式\eqref{align:ch05-5.4}是严格的。类似现象在函数序列中已遇到过（见\cref{example:ch02-弱收敛振荡}(i)）。
\end{remark}

\begin{remark}\label{remark:ch05-函数转测度}
设 $X$ 为局部紧 Hausdorff 空间，$\mu\in\mathscr{R}^+_f(X)$。设 $\{h_k\}\subseteq L^1(X,\mathscr{B}(X),\mu)$，对任意 $k\in\mathbb{N}$ 置
\begin{align}
\nu_k(E)\triangleq \int_E h_k\,\mathrm{d}\mu\qquad(E\in\mathscr{B}(X)).\label{align:ch05-5.5}
\end{align}
显然 $\{\nu_k\}\subseteq\mathscr{R}_f(X)$，且对所有 $k\in\mathbb{N}$ 有 $\nu_k\ll\mu$。$\{h_k\}$ 在 $L^1(X,\mathscr{B}(X),\mu)$ 中的收敛与 $\{\nu_k\}$ 在 $\mathscr{R}_f(X)$ 中的收敛之间存在明显联系：
\begin{itemize}
\item 若 $h_k\to h$ 在 $L^1(X)$ 中强收敛，且 $\nu(E)\triangleq \int_E h\,\mathrm{d}\mu$（$E\in\mathscr{B}(X)$），则 $\|\nu_k-\nu\|_{\mathscr{R}_f(X)}=|\nu_k-\nu|(X)=\int_X|h_k-h|\,\mathrm{d}\mu\to0$，故 $\nu_k\to\nu$ 在 $\mathscr{R}_f(X)$ 中强收敛；
\item 若 $h_k\to h$ 在 $L^1(X)$ 中弱收敛，则特别地
\begin{align*}
\lim_{k\to\infty}\int_X g h_k\,\mathrm{d}\mu=\int_X g h\,\mathrm{d}\mu
\quad\text{对任意 }g\in C_b(X),
\end{align*}
从而
\begin{align*}
\lim_{k\to\infty}\int_X g\,\mathrm{d}\nu_k=\int_X g\,\mathrm{d}\nu
\quad\text{对任意 }g\in C_b(X)
\end{align*}
其中 $\nu$ 同上。故 $\nu_k\stackrel{n}{\rightharpoonup}\nu$。
\end{itemize}
更一般地，即使序列 $\{h_k\}$ 在 $L^1(X)$ 中不收敛，由\eqref{align:ch05-5.5}定义的序列 $\{\nu_k\}$ 仍可（弱$^*$或窄）收敛于某个 $\nu\in\mathscr{R}_f(X)$（见\cref{example:ch05-测度逼近}(i)）。在此类情形中，按记号滥用，我们偶尔写作 $h_k\stackrel{*}{\rightharpoonup}\nu$ 或 $h_k\stackrel{n}{\rightharpoonup}\nu$。
\end{remark}

牢记\cref{remark:ch05-函数转测度}，我们指出一些与 $\mathscr{R}_f(X)$ 中收敛相关的现象。

\begin{example}\label{example:ch05-测度逼近}
\label{example:ch05-振荡集中}
(i) 每个 $\nu\in\mathscr{R}^+_f(\mathbb{R})$ 都是序列 $\{h_k\}\subseteq C^\infty_c(\mathbb{R})$（在 $L^1(\mathbb{R})\equiv L^1(\mathbb{R},\mathscr{B},\lambda)$ 中有界）的弱$^*$极限（更确切地说，是\eqref{align:ch05-5.5}中定义的密度 $h_k$ 所对应的测度序列 $\nu_k$ 的弱$^*$极限）。事实上，设 $\rho_{1/k}$ 为磨光子，$\rho_{1/k}(x)\triangleq k\rho(kx)$，其中 $\rho\in C^\infty_c(\mathbb{R})$，$\rho\geqslant0$，$\int_\mathbb{R}\rho\,\mathrm{d}\lambda=1$（$k\in\mathbb{N}$）。对任意 $k\in\mathbb{N}$，考虑卷积 $\nu*\rho_{1/k}\in C^\infty_c(\mathbb{R})$，
\begin{align*}
\int_\mathbb{R}g(x)(\nu*\rho_{1/k})(x)\,\mathrm{d}\lambda(x)
\triangleq \int_\mathbb{R}\mathrm{d}\nu(x)\int_\mathbb{R}\rho_{1/k}(y-x)g(y)\,\mathrm{d}\lambda(y)
\qquad(g\in C^\infty_c(\mathbb{R}))
\end{align*}
（见\cref{definition:ch03-磨光}(ii)与\cref{remark:ch03-磨光注}）。置 $h_k\triangleq \nu*\rho_{1/k}$，$\nu_k(E)\triangleq \int_E h_k\,\mathrm{d}\lambda$（$E\in\mathscr{B}$）。由\cref{proposition:ch03-磨光性质}(i)，
\begin{align*}
\lim_{k\to\infty}\int_\mathbb{R}g\,\mathrm{d}\nu_k
=\lim_{k\to\infty}\int_\mathbb{R}g h_k\,\mathrm{d}\lambda
=\int_\mathbb{R}g\,\mathrm{d}\nu
\quad\text{对任意 }g\in C^\infty_c(\mathbb{R}).
\end{align*}
易见 $\|h_k\|_{L^1(\mathbb{R})}\leqslant\|\nu\|=\nu(X)<\infty$，故由标准论证上述等式对任意 $g\in C_0(\mathbb{R})$ 成立。

(ii) 特别地，上述考虑适用于 $\nu$ 为原子测度的情形，例如 $\nu=\delta_0$。因此序列 $\{\nu_k\}\subseteq\mathscr{R}^+_f(\mathbb{R})$（对所有 $k\in\mathbb{N}$，$\nu_k\ll\lambda$）可弱$^*$收敛于测度 $\nu\perp\lambda$。在此情形我们也说发生了质量集中（见\cref{example:ch02-弱收敛振荡}）。反之，考虑测度空间 $(I,\mathscr{B}\cap I,\lambda|_{\mathscr{B}\cap I})$，其中 $I\equiv(0,1)$。对任意固定的 $h\in C_b(I)$，置 $\nu(E)\triangleq \int_E h\,\mathrm{d}\lambda$（$E\in\mathscr{B}\cap I$），并考虑原子测度序列
\begin{align*}
\nu_k\triangleq \frac{1}{k}\sum_{j=1}^k h\!\left(\frac{j}{k}\right)\delta_{j/k}\qquad(k\in\mathbb{N}).
\end{align*}
由初等结果，
\begin{align*}
\lim_{k\to\infty}\int_I g\,\mathrm{d}\nu_k
=\lim_{k\to\infty}\frac{1}{k}\sum_{j=1}^k g\!\left(\frac{j}{k}\right)h\!\left(\frac{j}{k}\right)
=\int_I gh\,\mathrm{d}\lambda=\int_I g\,\mathrm{d}\nu
\quad\text{对任意 }g\in C_b(I).
\end{align*}
故在此情形，对所有 $k\in\mathbb{N}$ 有 $\nu_k\perp\lambda$，$\nu\ll\lambda$，且 $\nu_k\stackrel{n}{\rightharpoonup}\nu$。
\end{example}

此后我们主要处理有限（正）Radon 测度序列。有限带号 Radon 测度序列的结果可由下文所证结果容易地得到（见\cref{remark:ch02-带号约化}）。

窄收敛可如下刻画。

\begin{theorem}[Portmanteau 定理]\label{theorem:ch05-portmanteau}
设 $X$ 为具可数基的局部紧 Hausdorff 空间，$\{\mu_k\}\subseteq\mathscr{R}^+_f(X)$，$\mu\in\mathscr{R}^+_f(X)$。则下列命题等价：
\begin{enumerate}[(i)]
\item\label{item:ch05-port-i} $\mu_k\stackrel{n}{\rightharpoonup}\mu$；
\item\label{item:ch05-port-ii} 对每个一致连续的 $g\in C_b(X)$，
\begin{align}
\lim_{k\to\infty}\int_X g\,\mathrm{d}\mu_k=\int_X g\,\mathrm{d}\mu;\label{align:ch05-5.6}
\end{align}
\item\label{item:ch05-port-iii} $\mu_k(X)\to\mu(X)$，且对任意闭集 $C\subseteq X$，
\begin{align}
\limsup_{k\to\infty}\mu_k(C)\leqslant\mu(C);\label{align:ch05-5.7}
\end{align}
\item\label{item:ch05-port-iv} $\mu_k(X)\to\mu(X)$，且对任意开集 $A\subseteq X$，
\begin{align}
\liminf_{k\to\infty}\mu_k(A)\geqslant\mu(A);\label{align:ch05-5.8}
\end{align}
\item\label{item:ch05-port-v} 对所有满足 $\mu(\partial E)=0$ 的 $E\in\mathscr{B}(X)$，
\begin{align}
\lim_{k\to\infty}\mu_k(E)=\mu(E).\label{align:ch05-5.9}
\end{align}
\end{enumerate}
\end{theorem}

\begin{proof}
(i)$\Rightarrow$(ii)。由\cref{definition:ch05-窄收敛}显然。

(ii)$\Rightarrow$(iii)。取 $g=1$ 表明 $\mu_k(X)\to\mu(X)$。设 $X$ 赋以与给定拓扑相容的度量 $d$。若 $C=\varnothing $，不等式\eqref{align:ch05-5.7}显然成立。否则，对任意 $n\in\mathbb{N}$，置 $C_{1/n}\triangleq \{x\in X\mid d(x,C)<1/n\}$。显然 $C_{1/n}$ 为开集，$C=\bigcap_{n=1}^\infty C_{1/n}$，故 $\mu(C)=\lim_{n\to\infty}\mu(C_{1/n})$。置 $f_n:X\to[0,\infty)$，$f_n(x)\triangleq \max\{1-nd(x,C),0\}$（$n\in\mathbb{N}$）。显然，对任意 $n\in\mathbb{N}$，函数 $f_n$ 在 $X$ 上一致连续，且 $\chi_C\leqslant f_n\leqslant\chi_{C_{1/n}}$。于是由\eqref{align:ch05-5.6}，对任意固定的 $n\in\mathbb{N}$，
\begin{align*}
\limsup_{k\to\infty}\mu_k(C)
\leqslant\lim_{k\to\infty}\int_X f_n\,\mathrm{d}\mu_k
=\int_X f_n\,\mathrm{d}\mu\leqslant\mu(C_{1/n}).
\end{align*}
令 $n\to\infty$ 即得\eqref{align:ch05-5.7}，结论得证。

(iii)$\Leftrightarrow$(iv)。由\eqref{align:ch05-5.7}，对任意开集 $A\subseteq X$，
\begin{align*}
\liminf_{k\to\infty}\mu_k(A)
=\lim_{k\to\infty}\mu_k(X)-\limsup_{k\to\infty}\mu_k(A^c)
\geqslant\mu(X)-\mu(A^c)=\mu(A),
\end{align*}
故\eqref{align:ch05-5.8}成立。反转论证即得逆命题。

(iv)$\Rightarrow$(v)。由假设 $\mu(\partial E)=\mu(\overline{E})-\mu(E^\circ)=0$，故 $\mu(E)=\mu(E^\circ)=\mu(\overline{E})$。于是由\eqref{align:ch05-5.7}–\eqref{align:ch05-5.8}，
\begin{align*}
\mu(E)=\mu(E^\circ)
\leqslant\liminf_{k\to\infty}\mu_k(E^\circ)
\leqslant\liminf_{k\to\infty}\mu_k(E)
\leqslant\limsup_{k\to\infty}\mu_k(E)
\leqslant\limsup_{k\to\infty}\mu_k(\overline{E})
\leqslant\mu(\overline{E})=\mu(E).
\end{align*}
故\eqref{align:ch05-5.9}成立。

(v)$\Rightarrow$(i)。对任意 $g\in C_b(X)$ 与 $t>0$，置 $E_t\triangleq \{g>t\}$。由于 $g$ 连续，对 $\lambda$-几乎所有 $t\in(0,\infty)$，有 $\mu(\partial E_t)\leqslant\mu(\{g=t\})=0$，故由\eqref{align:ch05-5.9}，
\begin{align}
\lim_{k\to\infty}\mu_k(\{g>t\})=\mu(\{g>t\})
\quad\text{对 }\lambda\text{-几乎所有 }t\in(0,\infty).\label{align:ch05-5.10}
\end{align}
于是由等式\eqref{align:ch02-2.19}（取 $\alpha=1$）、\eqref{align:ch05-5.10}及控制收敛定理，对任意 $g\in C_b(X)$ 有
\begin{align*}
\lim_{k\to\infty}\int_X g\,\mathrm{d}\mu_k
=\lim_{k\to\infty}\int_{[0,\infty)}\mu_k(\{g>t\})\,\mathrm{d}\lambda(t)
=\int_{[0,\infty)}\mu(\{g>t\})\,\mathrm{d}\lambda(t)
=\int_X g\,\mathrm{d}\mu.
\end{align*}
证明完成。
\end{proof}

\begin{remark}\label{remark:ch05-窄收敛范数}
如\cref{remark:ch05-范数下半连续}已述，\cref{theorem:ch05-portmanteau}表明：若 $\{\mu_k\}\subseteq\mathscr{R}^+_f(X)$，$\mu\in\mathscr{R}^+_f(X)$，且 $\mu_k\stackrel{n}{\rightharpoonup}\mu$，则 $\|\mu_k\|_{\mathscr{R}_f(X)}\to\|\mu\|_{\mathscr{R}_f(X)}$。
\end{remark}

\subsection{Prokhorov 距离}

本小节结果对度量空间 $(X,d)$ 表述。然而由\cref{proposition:appA-可度量化}，当 $X$ 为具可数基的局部紧 Hausdorff 空间且 $d$ 为与其拓扑相容的度量时，这些结果同样成立。

设 $(X,d)$ 为度量空间。对任意 $E\subseteq X$ 与 $r>0$，置 $E_r\triangleq \{x\in X\mid d(x,E)<r\}$。对任意 $\mu_1,\mu_2\in\mathscr{R}^+_f(X)$，定义
\begin{align}
d_P(\mu_1,\mu_2)\triangleq \inf\left\{r>0\ \middle|\
\mu_1(E)\leqslant\mu_2(E_r)+r,\ \mu_2(E)\leqslant\mu_1(E_r)+r
\ \forall E\in\mathscr{B}(X)\right\}.\label{align:ch05-5.11}
\end{align}
注意每个 $r\geqslant\max\{\|\mu_1\|_{\mathscr{R}_f(X)},\|\mu_2\|_{\mathscr{R}_f(X)}\}$ 都属于右端集合，故 $d_P(\mu_1,\mu_2)<\infty$。

\begin{proposition}\label{proposition:ch05-Prokhorov度量}
$(\mathscr{R}^+_f(X),d_P)$ 为度量空间。
\end{proposition}

\begin{proof}
我们证明映射 $(\mu_1,\mu_2)\mapsto d_P(\mu_1,\mu_2)$ 是 $\mathscr{R}^+_f(X)$ 上的度量。显然 $d_P(\mu_1,\mu_2)=d_P(\mu_2,\mu_1)$。此外：
\begin{enumerate}[(i)]
\item 对所有 $E\in\mathscr{B}(X)$，有 $E\subseteq E_r$，故对所有 $\mu\in\mathscr{R}^+_f(X)$ 与 $r>0$ 有 $\mu(E)\leqslant\mu(E_r)$。因此
\begin{align*}
d_P(\mu,\mu)\triangleq \inf\{r>0\mid\mu(E)\leqslant\mu(E_r)+r\ \forall E\in\mathscr{B}(X)\}=0;
\end{align*}
\item 若 $d_P(\mu_1,\mu_2)=0$，则存在序列 $r_n\to0^+$ 使对所有 $n$ 与 $E\in\mathscr{B}(X)$ 有 $\mu_1(E)\leqslant\mu_2(E_{r_n})+r_n$，$\mu_2(E)\leqslant\mu_1(E_{r_n})+r_n$。由于 $\bigcap_{n=1}^\infty\overline{E_{r_n}}=\overline{E}$，故对所有 $E\in\mathscr{B}(X)$ 有 $\mu_1(E)\leqslant\mu_2(E)$ 与 $\mu_2(E)\leqslant\mu_1(E)$，从而对所有闭集 $C\in\mathscr{B}(X)$ 有 $\mu_1(C)=\mu_2(C)$。特别地 $\mu_1(X)=\mu_2(X)$，故对任意开集 $A\in\mathscr{B}(X)$ 有 $\mu_1(A)=\mu_2(A)$。于是由 $\mu_1,\mu_2$ 的正则性（见\cref{lemma:ch01-正则性}(i)）得 $\mu_1=\mu_2$；
\item 设 $\mu_1,\mu_2,\mu_3\in\mathscr{R}^+_f(X)$，$r,s>0$ 使 $d_P(\mu_1,\mu_2)<r$，$d_P(\mu_2,\mu_3)<s$。则对所有 $E\in\mathscr{B}(X)$，
\begin{align*}
\mu_1(E)\leqslant\mu_2(E_r)+r\leqslant\mu_3(E_{r+s})+r+s,\qquad
\mu_3(E)\leqslant\mu_2(E_s)+s\leqslant\mu_1(E_{r+s})+r+s,
\end{align*}
故 $d_P(\mu_1,\mu_3)\leqslant r+s$。对 $r,s$ 取inf即得
\begin{align*}
d_P(\mu_1,\mu_3)\leqslant d_P(\mu_1,\mu_2)+d_P(\mu_2,\mu_3),
\end{align*}
即三角不等式。结论得证。
\end{enumerate}
\end{proof}

\begin{definition}\label{definition:ch05-Prokhorov}
\eqref{align:ch05-5.11}中定义的 $\mathscr{R}^+_f(X)$ 上的度量 $d_P$ 称为 \textbf{Prokhorov 度量}。
\end{definition}

我们现在证明如下结果。

\begin{lemma}\label{lemma:ch05-Prokhorov对称}
设 $(X,d)$ 为度量空间，$\mu_1,\mu_2\in\mathscr{R}^+_f(X)$ 满足对所有 $E\in\mathscr{B}(X)$，
\begin{align}
\mu_1(E)\leqslant\mu_2(E_r)+r\qquad(r>0).\label{align:ch05-5.12}
\end{align}
则对所有 $F\in\mathscr{B}(X)$，
\begin{align}
\mu_2(F)\leqslant\mu_1(F_r)+r+\|\mu_2\|_{\mathscr{R}_f(X)}-\|\mu_1\|_{\mathscr{R}_f(X)}.\label{align:ch05-5.13}
\end{align}
\end{lemma}

\begin{remark}\label{remark:ch05-Prokhorov等范数}
由\cref{lemma:ch05-Prokhorov对称}，若\eqref{align:ch05-5.12}成立且 $\|\mu_1\|_{\mathscr{R}_f(X)}=\|\mu_2\|_{\mathscr{R}_f(X)}$，则
\begin{align}
d_P(\mu_1,\mu_2)
=\inf\{r>0\mid\mu_1(E)\leqslant\mu_2(E_r)+r\ \forall E\in\mathscr{B}(X)\}
=\inf\{r>0\mid\mu_2(E)\leqslant\mu_1(E_r)+r\ \forall E\in\mathscr{B}(X)\}.\label{align:ch05-5.14}
\end{align}
\end{remark}

\begin{proof}[证明\cref{lemma:ch05-Prokhorov}]
在\eqref{align:ch05-5.12}中取 $E=(F_r)^c$，立得
\begin{align*}
\mu_1(F_r)=\|\mu_1\|_{\mathscr{R}_f(X)}-\mu_1((F_r)^c)
=\|\mu_1\|_{\mathscr{R}_f(X)}-\mu_1(E)
\geqslant\|\mu_1\|_{\mathscr{R}_f(X)}-\mu_2(E_r)-r\\
=\|\mu_1\|_{\mathscr{R}_f(X)}-\|\mu_2\|_{\mathscr{R}_f(X)}
+\mu_2((E_r)^c)-r
\geqslant\|\mu_1\|_{\mathscr{R}_f(X)}-\|\mu_2\|_{\mathscr{R}_f(X)}
+\mu_2(F)-r.
\end{align*}
（注意 $E\subseteq(F_r)^c\Leftrightarrow F\subseteq(E_r)^c$。）
\end{proof}

\begin{example}\label{example:ch05-Delta距离}
设 $(X,d)$ 为度量空间。则对所有 $x_1,x_2\in X$，
\begin{align}
d_P(\delta_{x_1},\delta_{x_2})=\min\{1,d(x_1,x_2)\}.\label{align:ch05-5.15}
\end{align}
为证\eqref{align:ch05-5.15}，由\eqref{align:ch05-5.14}注意到
\begin{align*}
d_P(\delta_{x_1},\delta_{x_2})
=\inf\{r>0\mid\chi_E(x_1)\leqslant\chi_{E_r}(x_2)+r\ \forall E\in\mathscr{B}(X)\}.
\end{align*}
由于 $\chi_E(x_1)\leqslant1$，有 $d_P(\delta_{x_1},\delta_{x_2})\leqslant1$。此外，若 $r>d(x_1,x_2)$，则对所有 $E\in\mathscr{B}(X)$ 有 $\chi_E(x_1)\leqslant\chi_{E_r}(x_2)+r$：事实上，若 $x_1\in E$（否则不等式平凡）且 $r>d(x_1,x_2)$，则 $x_2\in E_r$，故 $\chi_{E_r}(x_2)=1$，不等式成立。综上 $d_P(\delta_{x_1},\delta_{x_2})\leqslant\min\{1,d(x_1,x_2)\}$。

为证反向不等式，取 $E=\{x_1\}$，故 $\chi_E(x_1)=1$。若 $d(x_1,x_2)\geqslant1$ 且 $r\in(0,1)$，则 $x_2\in(E_r)^c$，故 $\chi_{E_r}(x_2)=0$。于是
\begin{align}
\chi_E(x_1)=1>r=\chi_{E_r}(x_2)+r,\label{align:ch05-5.16}
\end{align}
故 $d_P(\delta_{x_1},\delta_{x_2})\geqslant1=\min\{1,d(x_1,x_2)\}$。类似地，若 $d(x_1,x_2)<1$ 且 $r\in(0,d(x_1,x_2)]$，则 $x_2\in(E_r)^c$，不等式\eqref{align:ch05-5.16}再次成立。故 $d_P(\delta_{x_1},\delta_{x_2})\geqslant d(x_1,x_2)=\min\{1,d(x_1,x_2)\}$。于是等式\eqref{align:ch05-5.15}成立。
\end{example}

\begin{remark}\label{remark:ch05-Delta等距}
由\cref{example:ch05-Delta距离}，映射 $X\to\mathscr{R}^+_f(X)$，$x\mapsto\delta_x$，是 $(X,\min\{1,d\})$ 到 $(\mathscr{R}^+_f(X),d_P)$ 的等距嵌入。注意 $\min\{1,d\}$ 是 $X$ 上与 $d$ 等价的度量。
\end{remark}

现在我们可以证明 $\mathscr{R}^+_f(X)$ 上的窄拓扑是可度量化的。

\begin{proposition}\label{proposition:ch05-窄拓扑可度量化}
设 $(X,d)$ 为可分度量空间。则 $\mathscr{R}^+_f(X)$ 上的窄拓扑与 Prokhorov 度量相伴的度量拓扑一致。
\end{proposition}

\begin{remark}\label{remark:ch05-Prokhorov推论}
由\cref{proposition:ch05-窄拓扑可度量化}，若 $(X,d)$ 为可分度量空间，则：
\begin{enumerate}[(i)]
\item\label{item:ch05-prok-i} $\mu_k\stackrel{n}{\rightharpoonup}\mu$ 当且仅当 $d_P(\mu_k,\mu)\to0$；
\item\label{item:ch05-prok-ii} 子集 $M\subseteq\mathscr{R}^+_f(X)$ 在窄拓扑中相对序列紧，当且仅当它在 $(\mathscr{R}^+_f(X),d_P)$ 中相对紧。
\end{enumerate}
\end{remark}

为证\cref{proposition:ch05-窄拓扑可度量化}，我们先建立如下引理。

\begin{lemma}\label{lemma:ch05-Prokhorov到窄}
设 $(X,d)$ 为度量空间，$\{\mu_k\}\subseteq\mathscr{R}^+_f(X)$，$\mu\in\mathscr{R}^+_f(X)$ 满足 $d_P(\mu_k,\mu)\to0$。则 $\mu_k\stackrel{n}{\rightharpoonup}\mu$。
\end{lemma}

\begin{proof}
设 $\{\epsilon_k\}\subseteq(0,\infty)$ 满足当 $k\to\infty$ 时 $\epsilon_k\to0^+$，且对所有 $k\in\mathbb{N}$ 有 $d_P(\mu_k,\mu)<\epsilon_k$。对任意闭集 $C\subseteq X$，置 $C_k\equiv C_{\epsilon_k}\triangleq \{x\in X\mid d(x,C)<\epsilon_k\}$，故 $C=\bigcap_{k=1}^\infty C_k$ 且 $\mu(C)=\lim_{k\to\infty}\mu(C_k)$。由于 $d_P(\mu_k,\mu)<\epsilon_k$，由定义\eqref{align:ch05-5.11}，对任意 $k\in\mathbb{N}$ 有
\begin{align}
\mu_k(C)\leqslant\mu(C_k)+\epsilon_k,\qquad
\mu(C)\leqslant\mu_k(C_k)+\epsilon_k.\label{align:ch05-5.17}
\end{align}
由于 $\lim_{k\to\infty}\mu(C_k)=\mu(C)$ 且 $\lim_{k\to\infty}\epsilon_k=0$，在\eqref{align:ch05-5.17}第一式中令 $k\to\infty$ 得\eqref{align:ch05-5.7}。另一方面，在\eqref{align:ch05-5.17}第二式与\eqref{align:ch05-5.7}中均取 $C=X$，得
\begin{align*}
\mu(X)\leqslant\liminf_{k\to\infty}\mu_k(X)\leqslant\limsup_{k\to\infty}\mu_k(X)\leqslant\mu(X),
\end{align*}
故 $\lim_{k\to\infty}\mu_k(X)=\mu(X)$。于是由\cref{theorem:ch05-portmanteau}（见断言 (i) 与 (iii)）结论得证。
\end{proof}

\begin{proof}[证明\cref{proposition:ch05-窄拓扑可度量化}]
由\cref{lemma:ch05-Prokhorov到窄}，我们只需证明 $\mu_k\stackrel{n}{\rightharpoonup}\mu$ 蕴含 $d_P(\mu_k,\mu)\to0$。为此，设 $D\equiv\{x_k\}$ 为 $X$ 的可数稠密子集。

对任意固定的 $\epsilon>0$，置 $E_1\triangleq B_{\epsilon/2}(x_1)$，$E_{m+1}\triangleq B_{\epsilon/2}(x_{m+1})\setminus\bigcup_{k=1}^m E_k$（$m\in\mathbb{N}$）。显然，对所有 $m\in\mathbb{N}$ 有 $\operatorname{diam}(E_m)\leqslant\epsilon$，对所有 $l,m\in\mathbb{N}$，$l\neq m$ 有 $E_l\cap E_m=\varnothing $，且 $\bigcup_{m=1}^\infty E_m=X$。由于 $\mu(X)=\sum_{m=1}^\infty\mu(E_m)<\infty$，存在 $\bar{m}\in\mathbb{N}$ 使
\begin{align}
\mu\!\left(\bigcup_{m=\bar{m}+1}^\infty E_m\right)
=\sum_{m=\bar{m}+1}^\infty\mu(E_m)<\epsilon.\label{align:ch05-5.18}
\end{align}
考虑有限族
\begin{align*}
S\triangleq \left\{U=E_{m_1}\cup E_{m_2}\cup\cdots\cup E_{m_p}\ \middle|\
1\leqslant m_1<m_2<\cdots<m_p\leqslant\bar{m}\right\}\subseteq\mathcal{P}(X),
\end{align*}
对任意 $U\in S$，置 $U_\epsilon\triangleq \{x\in X\mid d(x,U)<\epsilon\}$。由于 $U_\epsilon$ 为开集且由假设 $\mu_k\stackrel{n}{\rightharpoonup}\mu$，由 Portmanteau 定理 $\liminf_{k\to\infty}\mu_k(U_\epsilon)\geqslant\mu(U_\epsilon)$。因此，由于 $S$ 有限，存在 $\bar{k}\in\mathbb{N}$ 使
\begin{align}
\mu_k(U_\epsilon)\geqslant\mu(U_\epsilon)-\epsilon
\quad\text{对所有 }k\geqslant\bar{k}\text{ 与 }U\in S.\label{align:ch05-5.19}
\end{align}
现注意到对所有 $F\in\mathscr{B}(X)$，
\begin{align}
F=\left[F\cap\bigcup_{m=1}^{\bar{m}}E_m\right]
\cup\left[F\cap\bigcup_{m=\bar{m}+1}^\infty E_m\right]
\subseteq[F\cap U_\epsilon]\cup\left[F\cap\bigcup_{m=\bar{m}+1}^\infty E_m\right]\label{align:ch05-5.20}
\end{align}
对某个 $U=E_{m_1}\cup E_{m_2}\cup\cdots\cup E_{m_p}\in S$（使对所有 $j=m_1,\cdots ,m_p$ 有 $E_{m_j}\cap F\neq\varnothing $）成立（注意这蕴含 $U_\epsilon\subseteq F_{2\epsilon}$，因为对任意 $m\in\mathbb{N}$ 有 $\operatorname{diam}(E_m)\leqslant\epsilon$）。由\eqref{align:ch05-5.18}–\eqref{align:ch05-5.20}得
\begin{align}
\mu(F)\leqslant\mu(U_\epsilon)+\epsilon
\leqslant\mu_k(U_\epsilon)+2\epsilon
\leqslant\mu_k(F_{2\epsilon})+2\epsilon
\quad\text{对所有 }k\geqslant\bar{k}.\label{align:ch05-5.21}
\end{align}
由 $F\in\mathscr{B}(X)$ 的任意性与\cref{lemma:ch05-Prokhorov对称}，从\eqref{align:ch05-5.21}得
\begin{align}
\mu_k(G)\leqslant\mu(G_{2\epsilon})+2\epsilon
+\|\mu_k\|_{\mathscr{R}_f(X)}-\|\mu\|_{\mathscr{R}_f(X)}
\leqslant\mu(G_{3\epsilon})+3\epsilon\label{align:ch05-5.22}
\end{align}
对所有 $G\in\mathscr{B}(X)$ 与所有充分大的 $k\in\mathbb{N}$ 成立，因为 $\lim_{k\to\infty}\|\mu_k\|_{\mathscr{R}_f(X)}=\|\mu\|_{\mathscr{R}_f(X)}$（见\cref{remark:ch05-窄收敛范数}）。

另一方面，由\eqref{align:ch05-5.21}，对所有充分大的 $k\in\mathbb{N}$ 与任意 $G\in\mathscr{B}(X)$，我们同样有
\begin{align}
\mu(G)\leqslant\mu_k(G_{3\epsilon})+3\epsilon.\label{align:ch05-5.23}
\end{align}
由定义\eqref{align:ch05-5.11}，从\eqref{align:ch05-5.22}–\eqref{align:ch05-5.23}可得：对所有 $\epsilon>0$ 与充分大的 $k\in\mathbb{N}$，$d_P(\mu_k,\mu)<\epsilon$。结论得证。
\end{proof}

\subsection{窄收敛与胎紧性}

刻画 $\mathscr{R}_f(X)$ 在弱$^*$与窄两种拓扑下的相对序列紧子集是有趣的。关于弱$^*$收敛，我们有如下结果：

\begin{theorem}\label{theorem:ch05-弱星紧性}
\label{theorem:ch05-弱星紧}
设 $X$ 为具可数基的局部紧 Hausdorff 空间，$M\subseteq\mathscr{R}_f(X)$。则下列命题等价：
\begin{enumerate}[(i)]
\item\label{item:ch05-wstar-i} $M$ 在弱$^*$拓扑中相对序列紧；
\item\label{item:ch05-wstar-ii} $M$ 有界。
\end{enumerate}
\end{theorem}

\begin{proof}
(i)$\Rightarrow$(ii)。由\cref{theorem:ch02-Rf对偶}，Banach 空间 $\mathscr{R}_f(X)$ 是 $C_0(X)$ 的对偶空间。若 $M$ 无界，则存在 $\{\mu_k\}\subseteq M$ 使 $\|\mu_k\|_{\mathscr{R}_f(X)}\to\infty$。另一方面，由假设，每个 $\{\mu_k\}\subseteq M$ 都含一个在 $\mathscr{R}_f(X)$ 中弱$^*$收敛的子列，因而由一致有界原理有界。矛盾证明了断言。

(ii)$\Rightarrow$(i) 由 Banach–Alaoglu 定理得到。
\end{proof}

关于窄收敛，我们从如下定义开始：

\begin{definition}\label{definition:ch05-胎紧}
设 $X$ 为局部紧 Hausdorff 空间。子集 $M\subseteq\mathscr{R}_f(X)$ 称为\textbf{胎紧的}（tight），若对任意 $\epsilon>0$，存在紧子集 $K\subseteq X$ 使对所有 $\mu\in M$ 有 $|\mu|(K^c)<\epsilon$。
\end{definition}

\begin{example}\label{example:ch05-胎紧}
设 $X=\mathbb{R}$，$\mu=\delta_{\bar{x}}$（$\bar{x}\in\mathbb{R}$）。故对所有 $E\in\mathscr{B}(\mathbb{R})$ 有 $\mu(E)=\chi_E(\bar{x})$。显然，对任意 $F\subseteq\mathbb{R}$，集合 $M=\{\delta_{\bar{x}}\mid\bar{x}\in F\}$ 胎紧当且仅当 $F$ 有界。注意若 $F$ 有界，则集合 $\{\alpha\delta_{\bar{x}}\mid\bar{x}\in F,\ \alpha>0\}$ 胎紧但无界。
\end{example}

为后续参考，我们证明如下结果。

\begin{proposition}\label{proposition:ch05-胎紧拓扑一致}
设 $X$ 为具可数基的局部紧 Hausdorff 空间，$M\subseteq\mathscr{R}^+_f(X)$ 胎紧。则 $M$ 上的窄拓扑与 $M$ 上的弱$^*$拓扑一致，即 $\mathcal{T}_n\cap M=\mathcal{T}_{w^*}\cap M$。
\end{proposition}

\begin{proof}
只需证明：对任意序列 $\{\mu_k\}\subseteq M$，若 $\mu_k\stackrel{*}{\rightharpoonup}\mu$，则 $\mu_k\stackrel{n}{\rightharpoonup}\mu$。固定 $\epsilon>0$，由\cref{definition:ch05-胎紧}存在紧子集 $K\subseteq X$ 使对所有 $k\in\mathbb{N}$ 有 $\mu_k(K^c)<\epsilon$。另一方面，由\cref{proposition:appA-Urysohn}，对任意开集 $A\subseteq X$，$A\supseteq K$，存在 $f\in C_c(X)$ 使 $f(X)\subseteq[0,1]$，$f|_K=1$，且 $\operatorname{supp}f\subseteq A$。固定 $g\in C_b(X)$。则 $fg\in C_0(X)$，且
\begin{align*}
\left|\int_X g\,\mathrm{d}\mu_k-\int_X g\,\mathrm{d}\mu\right|
\leqslant\left|\int_K g\,\mathrm{d}\mu_k-\int_K g\,\mathrm{d}\mu\right|
+\left|\int_{K^c}g\,\mathrm{d}\mu_k-\int_{K^c}g\,\mathrm{d}\mu\right|\\
\leqslant\left|\int_X fg\,\mathrm{d}\mu_k-\int_X fg\,\mathrm{d}\mu\right|
+\left|\int_X[\chi_K-f]g\,\mathrm{d}\mu_k\right|
+\left|\int_X[\chi_K-f]g\,\mathrm{d}\mu\right|
+\left|\int_{K^c}g\,\mathrm{d}\mu_k\right|
+\left|\int_{K^c}g\,\mathrm{d}\mu\right|
=:I_1+I_2+I_3+I_4+I_5.
\end{align*}
由于由假设 $\mu_k\stackrel{*}{\rightharpoonup}\mu$，存在 $\bar{k}\in\mathbb{N}$ 使对所有 $k>\bar{k}$，
\begin{align*}
I_1=\left|\int_X fg\,\mathrm{d}\mu_k-\int_X fg\,\mathrm{d}\mu\right|<\epsilon.
\end{align*}
我们还有，对所有 $k\in\mathbb{N}$，
\begin{align*}
I_2+I_4
=\left|\int_X[\chi_K-f]g\,\mathrm{d}\mu_k\right|
+\left|\int_{K^c}g\,\mathrm{d}\mu_k\right|
\leqslant 2\|g\|_\infty\mu_k(K^c)<2\|g\|_\infty\epsilon.
\end{align*}
此外，
\begin{align*}
I_3+I_5
=\left|\int_X[\chi_K-f]g\,\mathrm{d}\mu\right|
+\left|\int_{K^c}g\,\mathrm{d}\mu\right|
\leqslant 2\|g\|_\infty\mu(K^c)\leqslant 2\|g\|_\infty\epsilon,
\end{align*}
因为由范数的下半连续性（见\cref{remark:ch05-范数下半连续}），$\mu(K^c)\leqslant\liminf_{k\to\infty}\mu_k(K^c)$。

综上，我们证明了：对任意固定的 $g\in C_b(X)$ 与任意 $\epsilon>0$，存在 $\bar{k}\in\mathbb{N}$ 使对所有 $k>\bar{k}$，
\begin{align*}
\left|\int_X g\,\mathrm{d}\mu_k-\int_X g\,\mathrm{d}\mu\right|<(1+4\|g\|_\infty)\epsilon.
\end{align*}
结论得证。
\end{proof}

如下结果是\cref{theorem:ch05-弱星紧性}在弱$^*$收敛方面的对应物，表明窄收敛与胎紧性紧密交织。

\begin{theorem}[Prokhorov]\label{theorem:ch05-prokhorov}\label{lemma:ch05-Prokhorov}
设 $X$ 为具可数基的局部紧 Hausdorff 空间，$M\subseteq\mathscr{R}^+_f(X)$。则下列命题等价：
\begin{enumerate}[(i)]
\item\label{item:ch05-prokth-i} $M$ 在窄拓扑中相对序列紧；
\item\label{item:ch05-prokth-ii} $M$ 有界且胎紧。
\end{enumerate}
\end{theorem}

\begin{remark}\label{remark:ch05-polish}
更一般地，若 $X$ 为 Polish 空间（见附录 A,\cref{subsection:appA-Polish}），\cref{theorem:ch05-prokhorov}仍成立。\cref{proposition:ch05-ulam}亦然。
\end{remark}

\cref{theorem:ch05-prokhorov}的证明分两步给出，对应于\cref{proposition:ch05-ulam}与\cref{proposition:ch05-有界胎紧紧}。

\begin{proposition}[Ulam]\label{proposition:ch05-ulam}
设 $X$ 为具可数基的局部紧 Hausdorff 空间，$M\subseteq\mathscr{R}^+_f(X)$ 在窄拓扑中相对序列紧。则 $M$ 有界且胎紧。
\end{proposition}

\begin{proof}
由假设，对任意序列 $\{\mu_k\}\subseteq M$，存在 $\{\mu_{k_l}\}\subseteq\{\mu_k\}$ 与 $\mu\in\mathscr{R}_f(X)$ 使 $\mu_{k_l}\stackrel{n}{\rightharpoonup}\mu$，因而 $\|\mu_{k_l}\|_{\mathscr{R}_f(X)}\to\|\mu\|_{\mathscr{R}_f(X)}$（见\cref{remark:ch05-窄收敛范数}）。显然这蕴含 $M$ 有界。

为证 $M$ 胎紧，初步注意 $X$ 为 Polish 空间，故存在相容度量 $d$ 使 $(X,d)$ 完备。固定 $\epsilon>0$，设 $D\equiv\{x_k\}$ 为 $X$ 的可数稠密子集。对任意固定的 $p\in\mathbb{N}$，置 $E_{n,p}\triangleq \bigcup_{k=1}^n B_{1/p}(x_k)$（其中 $B_{1/p}(x_k)\triangleq \{x\in X\mid d(x_k,x)<1/p\}$；$n\in\mathbb{N}$）。显然每个 $E_{n,p}$ 为开集，对所有 $n\in\mathbb{N}$ 有 $E_{n,p}\subseteq E_{n+1,p}$，且 $\bigcup_{n=1}^\infty E_{n,p}=X$。我们证明如下断言。

\textbf{断言。} 设 $\{A_n\}\subseteq\mathcal{P}(X)$ 为递增开集序列使 $\bigcup_{n=1}^\infty A_n=X$。则对每个 $\epsilon>0$，存在 $\bar{n}\in\mathbb{N}$ 使对所有 $\mu\in M$ 有 $\mu(A_{\bar{n}}^c)<\epsilon$。

由断言，对任意 $p\in\mathbb{N}$，存在 $n_p\in\mathbb{N}$ 使对所有 $\mu\in M$ 有 $\mu((E_{n_p,p})^c)<\epsilon/2^p$，特别地
\begin{align}
\mu((E_{n_p,p})^c)<\frac{\epsilon}{2^p}
\quad\text{对所有 }\mu\in M.\label{align:ch05-5.24}
\end{align}
置
\begin{align*}
K\triangleq \bigcap_{p=1}^\infty E_{n_p,p}\subseteq\bigcap_{p=1}^\infty\bigcup_{k=1}^{n_p}B_{1/p}(x_k).
\end{align*}
由\eqref{align:ch05-5.24}，对所有 $\mu\in M$ 有 $\mu(K^c)<\epsilon$，故若我们证明 $K$ 紧，结论即得。为此注意，对每个 $p\in\mathbb{N}$，集合 $K$ 包含在有限并 $\bigcup_{k=1}^{n_p}B_{1/p}(x_k)$ 中，且 $\operatorname{diam}B_{1/p}(x_k)=2/p$（$p\in\mathbb{N}$）。由对角论证，每个序列 $\{x'_j\}\subseteq K$ 含 Cauchy 子列 $\{x'_{j_l}\}\subseteq\{x'_j\}$。由于 $K$ 闭且 $X$ 完备，存在 $x'\in K$ 使当 $l\to\infty$ 时 $d(x'_{j_l},x')\to0$。故 $K$ 序列紧，从而紧。结论得证。

尚需证明断言。反设存在递增开集序列 $\{A_n\}$ 使 $\bigcup_{n=1}^\infty A_n=X$，及 $\epsilon>0$ 具有如下性质：对所有 $n\in\mathbb{N}$，存在 $\mu_n\in M$ 使 $\mu_n(A_n^c)\geqslant\epsilon$。由于 $M$ 在窄拓扑中相对序列紧，存在 $\{\mu_{n_k}\}\subseteq\{\mu_n\}$ 与 $\mu\in\mathscr{R}_f(X)$ 使 $\mu_{n_k}\stackrel{n}{\rightharpoonup}\mu$。于是由\cref{theorem:ch05-portmanteau}与 $\{A_n\}$ 的单调性，对每个固定的 $n\in\mathbb{N}$，
\begin{align*}
\mu(A_n^c)\geqslant\limsup_{k\to\infty}\mu_{n_k}(A_n^c)
\geqslant\limsup_{k\to\infty}\mu_{n_k}(A_{n_k}^c)\geqslant\epsilon.
\end{align*}
然而由于 $\mu$ 有限，我们得 $\lim_{n\to\infty}\mu(A_n^c)=0$，矛盾。故断言成立，证明完成。
\end{proof}

\begin{proposition}\label{proposition:ch05-有界胎紧紧}
设 $X$ 为具可数基的局部紧 Hausdorff 空间，$M\subseteq\mathscr{R}^+_f(X)$ 有界且胎紧。则 $M$ 在窄拓扑中相对序列紧。
\end{proposition}

\begin{proof}
由假设，
\begin{align}
\sup_{\mu\in M}\|\mu\|_{\mathscr{R}_f(X)}
=\sup_{\mu\in M}\mu(X)=:C<\infty.\label{align:ch05-5.25}
\end{align}
回忆由\cref{proposition:appA-紧生成}，空间 $X$ 为 $\sigma$-紧。我们先在 $X$ 紧的假设下证明结果。设 $D\equiv\{f_k\}$ 为 $C_b(X)$ 的可数稠密子集（由于 $X$ 紧，$C_b(X)$ 可分；见\cref{proposition:appA-紧可分}）。则
\begin{align*}
\sup_{\mu\in M}\left|\int_X f_1\,\mathrm{d}\mu\right|
\leqslant\|f_1\|_\infty\sup_{\mu\in M}\|\mu\|_{\mathscr{R}_f(X)}
=\|f_1\|_\infty C<\infty,
\end{align*}
故集合 $\{\int_X f_1\,\mathrm{d}\mu\mid\mu\in M\}\subseteq\mathbb{R}$ 有界。于是存在序列 $\{\mu_j^1\}\subseteq M$ 与 $a_1\in\mathbb{R}$ 使
\begin{align}
\lim_{j\to\infty}\int_X f_1\,\mathrm{d}\mu_j^1=a_1.\label{align:ch05-5.26}
\end{align}
由于
\begin{align*}
\sup_{j\in\mathbb{N}}\left|\int_X f_2\,\mathrm{d}\mu_j^1\right|
\leqslant\|f_2\|_\infty\sup_{j\in\mathbb{N}}\|\mu_j^1\|
\leqslant\|f_2\|_\infty C<\infty,
\end{align*}
集合 $\{\int_X f_2\,\mathrm{d}\mu_j^1\mid j\in\mathbb{N}\}\subseteq\mathbb{R}$ 也有界。于是存在子列 $\{\mu_j^2\}\subseteq\{\mu_j^1\}$ 与 $a_2\in\mathbb{R}$ 使
\begin{align*}
\lim_{j\to\infty}\int_X f_2\,\mathrm{d}\mu_j^2=a_2.
\end{align*}
显然由\eqref{align:ch05-5.26}，我们也有 $\lim_{j\to\infty}\int_X f_1\,\mathrm{d}\mu_j^2=a_1$。迭代此论证，对每个 $k\in\mathbb{N}$，$k\geqslant3$，存在子列 $\{\mu_j^k\}\subseteq\{\mu_j^{k-1}\}$ 与 $a_k\in\mathbb{R}$ 使
\begin{align*}
\lim_{j\to\infty}\int_X f_l\,\mathrm{d}\mu_j^k=a_l
\quad\text{对所有 }l=1,\cdots ,k.
\end{align*}
由对角论证，序列 $\{\mu_j\}\subseteq M$（其中 $\mu_j\triangleq \mu_j^j$，$j\in\mathbb{N}$）满足
\begin{align}
\lim_{j\to\infty}\int_X f_k\,\mathrm{d}\mu_j=a_k
\quad\text{对所有 }k\in\mathbb{N}.\label{align:ch05-5.27}
\end{align}
现考虑映射 $F_0:D\to\mathbb{R}$，$F_0(f_k)\triangleq a_k$（$k\in\mathbb{N}$）。由\eqref{align:ch05-5.27}，$F_0$ 线性、正且连续，因为
\begin{align*}
|F_0(f_k)|=|a_k|
\leqslant\lim_{j\to\infty}\left|\int_X f_k\,\mathrm{d}\mu_j\right|
\leqslant\|f_k\|_\infty C
\quad\text{对所有 }k\in\mathbb{N}.
\end{align*}
由于 $D$ 在 $C_b(X)$ 中稠密，$F_0$ 可唯一延拓为 $(C_b(X))^*$ 上的泛函 $F$，它是正的且满足 $\|F\|\leqslant C$。我们证明如下断言。

\textbf{断言。} 对任意 $\epsilon>0$，存在紧子集 $K\subseteq X$ 使对所有 $f\in C_b(X)$ 满足 $f(X)\subseteq[0,1]$ 且 $f|_K=0$，有 $\langle F,f\rangle<\epsilon$。

依靠断言，我们可在 $X$ 紧时完成证明。事实上，由\cref{theorem:ch02-Riesz表示}，存在唯一 $\mu\in\mathscr{R}^+_f(X)$ 使
\begin{align}
\langle F,f\rangle=\int_X f\,\mathrm{d}\mu
\quad\text{对所有 }f\in C_b(X)\label{align:ch05-5.28}
\end{align}
（见\eqref{align:ch02-2.49}）；此外由\cref{theorem:ch02-Riesz范数}(ii)，$\|\mu\|_{\mathscr{R}_f(X)}=\|F\|\leqslant C$。设 $g\in C_b(X)$，$\{f_l\}\subseteq D$ 为使 $\lim_{l\to\infty}\|f_l-g\|_\infty=0$ 的序列。固定 $\epsilon>0$，取 $l_0\in\mathbb{N}$ 充分大使 $\|f_{l_0}-g\|_\infty<\epsilon/(4C)$。由\eqref{align:ch05-5.27}–\eqref{align:ch05-5.28}，
\begin{align*}
\langle F,f_{l_0}\rangle=F_0(f_{l_0})=a_{l_0}
=\int_X f_{l_0}\,\mathrm{d}\mu
=\lim_{j\to\infty}\int_X f_{l_0}\,\mathrm{d}\mu_j,
\end{align*}
故存在 $j_0\in\mathbb{N}$ 使
\begin{align*}
\left|\int_X f_{l_0}\,\mathrm{d}\mu_j-\int_X f_{l_0}\,\mathrm{d}\mu\right|<\frac{\epsilon}{2}
\quad\text{对所有 }j>j_0.
\end{align*}
则对所有 $j>j_0$，
\begin{align*}
\left|\int_X g\,\mathrm{d}\mu_j-\int_X g\,\mathrm{d}\mu\right|
\leqslant\left|\int_X(g-f_{l_0})\,\mathrm{d}\mu_j\right|
+\left|\int_X(g-f_{l_0})\,\mathrm{d}\mu\right|
+\left|\int_X f_{l_0}\,\mathrm{d}\mu_j-\int_X f_{l_0}\,\mathrm{d}\mu\right|\\
\leqslant\|f_{l_0}-g\|_\infty\left(\sup_{j\in\mathbb{N}}\|\mu_j\|+\|\mu\|_{\mathscr{R}_f(X)}\right)+\frac{\epsilon}{2}
<2C\cdot\frac{\epsilon}{4C}+\frac{\epsilon}{2}=\epsilon.
\end{align*}
综上，我们证明了存在序列 $\{\mu_j\}\subseteq M$ 与 $\mu\in\mathscr{R}^+_f(X)$ 使 $\mu_j\stackrel{n}{\rightharpoonup}\mu$，结论得证。

现证明上述断言。由于由假设 $M$ 胎紧，对任意 $\epsilon>0$，存在紧子集 $K\subseteq X$ 使 $\sup_{\mu\in M}\mu(K^c)<\epsilon/6$。设 $f\in C_b(X)$ 满足 $f(X)\subseteq[0,1]$ 且 $f|_K=0$。由于 $D$ 在 $C_b(X)$ 中稠密，存在序列 $\{f_l\}\subseteq D$（对所有 $l\in\mathbb{N}$，$f_l\geqslant0$）使 $\lim_{l\to\infty}\|f_l-f\|_\infty=0$；因而特别地，对所有充分大的 $l\in\mathbb{N}$，$\|f_l\|_\infty<1+\epsilon/6$，$\|f_l|_K\|_\infty<\epsilon/(6C)$。

则对任意充分大的 $l\in\mathbb{N}$，
\begin{align*}
\langle F,f_l\rangle=a_l
=\lim_{j\to\infty}\int_X f_l\,\mathrm{d}\mu_j
\leqslant\limsup_{j\to\infty}\int_K f_l\,\mathrm{d}\mu_j
+\limsup_{j\to\infty}\int_{K^c}f_l\,\mathrm{d}\mu_j\\
\leqslant\frac{\epsilon}{6C}\sup_{j\in\mathbb{N}}\mu_j(X)
+\left(1+\frac{\epsilon}{6}\right)\sup_{j\in\mathbb{N}}\mu_j(K^c)
<\frac{\epsilon}{6}+\left(1+\frac{\epsilon}{6}\right)\frac{\epsilon}{6}
<\frac{\epsilon}{2}.
\end{align*}
由于 $\|f_l-f\|_\infty\to0$ 且 $F$ 连续，由该不等式得
\begin{align*}
\langle F,f\rangle=\lim_{l\to\infty}\langle F,f_l\rangle
\leqslant\frac{\epsilon}{2}<\epsilon,
\end{align*}
故断言成立。$X$ 紧时的证明完成。

一般情形，设 $\{K_l\}$ 为 $X$ 的递增紧子集序列使 $\bigcup_{l=1}^\infty K_l=X$。由\eqref{align:ch05-5.25}，对每个 $l\in\mathbb{N}$，
\begin{align*}
\sup_{\mu\in M}\|\mu\|_{\mathscr{R}_f(K_l)}
\leqslant\sup_{\mu\in M}\|\mu\|_{\mathscr{R}_f(X)}=C.
\end{align*}
由上述考虑，对每个 $l\in\mathbb{N}$，存在序列 $\{\mu_j^l\}\subseteq M$ 与 $\mu^l\in\mathscr{R}^+_f(X)$ 使当 $j\to\infty$ 时 $\mu_j^l\stackrel{n}{\rightharpoonup}\mu^l$。于是由对角论证结论得证。
\end{proof}

由\cref{theorem:ch05-prokhorov}，拥有胎紧性判据是有用的。我们给出如下定义。

\begin{definition}\label{definition:ch05-下紧}
设 $X$ 为局部紧 Hausdorff 空间。函数 $\varphi:X\to(-\infty,\infty]$ 称为\textbf{下紧的}（inf-compact），若对每个 $r\in\mathbb{R}$，集合 $\{\varphi\leqslant r\}\triangleq \{x\in X\mid\varphi(x)\leqslant r\}$ 为紧集。
\end{definition}

\begin{example}\label{example:ch05-下紧}
若 $X=\mathbb{R}$，则每个连续函数 $\varphi:\mathbb{R}\to[0,\infty)$ 满足当 $|x|\to\infty$ 时 $\varphi(x)\to\infty$ 都是下紧的。事实上，对每个 $r\in\mathbb{R}$，集合 $\{\varphi\leqslant r\}$ 闭且有界，故紧。
\end{example}

\cref{definition:ch05-下紧}导出 $\mathscr{R}^+_f(X)$ 胎紧子集的如下刻画。

\begin{proposition}\label{proposition:ch05-胎紧判据}
设 $X$ 为具可数基的局部紧 Hausdorff 空间。则下列命题等价：
\begin{enumerate}[(i)]
\item\label{item:ch05-tight-i} $M\subseteq\mathscr{R}^+_f(X)$ 胎紧；
\item\label{item:ch05-tight-ii} 存在下紧函数 $\varphi:X\to[0,\infty)$ 使
\begin{align}
\sup_{\mu\in M}\int_X\varphi\,\mathrm{d}\mu<\infty.\label{align:ch05-5.29}
\end{align}
\end{enumerate}
\end{proposition}

\begin{proof}
(i)$\Rightarrow$(ii)。由\cref{proposition:appA-紧生成}，存在 $X$ 的递增紧子集序列 $\{K_n\}$ 使对所有 $\mu\in M$ 有 $\mu(K_n^c)<1/2^n$。置 $\varphi\triangleq \sum_{n=1}^\infty\chi_{K_n^c}$。则对每个 $\mu\in M$，
\begin{align*}
\sup_{\mu\in M}\int_X\varphi\,\mathrm{d}\mu
\leqslant\sum_{n=1}^\infty\frac{1}{2^n}<\infty,
\end{align*}
故不等式\eqref{align:ch05-5.29}成立。

为证 $\varphi$ 下紧，固定 $x\in X$。对每个 $n\in\mathbb{N}$，我们有 (a) $x\in K_n$ 或 (b) $x\in K_n^c$。由于每个 $n$ 有 $K_n^c\supseteq K_{n+1}^c$，
\begin{itemize}
\item 在情形 (a)，$\varphi(x)\leqslant n-1$；
\item 在情形 (b)，若 $x\in K_n^c\setminus K_{n+1}^c=K_n^c\cap K_{n+1}$ 则 $\varphi(x)=n$，若 $x\in K_{n+1}^c$ 则 $\varphi(x)\geqslant n+1$。
\end{itemize}
由此 $\{\varphi\leqslant n\}=K_{n+1}$，故显然对所有 $r>0$ 有 $\{\varphi\leqslant r\}=K_{[r]+1}$，其中 $[r]$ 表示小于 $r$ 的最大整数。由于 $K_{[r]+1}$ 紧，断言得证。

(ii)$\Rightarrow$(i)。由假设存在 $\varphi:X\to[0,\infty)$ 使对每个 $r>0$ 集合 $K_r\equiv\{\varphi\leqslant r\}$ 紧，且\eqref{align:ch05-5.29}成立。则由 Chebyshev 不等式，对每个 $\mu\in M$，
\begin{align*}
r\mu(K_r^c)<\int_{K_r^c}\varphi\,\mathrm{d}\mu
\leqslant\int_X\varphi\,\mathrm{d}\mu
\leqslant\sup_{\mu\in M}\int_X\varphi\,\mathrm{d}\mu,
\end{align*}
由 $r$ 的任意性结论得证。
\end{proof}

\begin{remark}\label{remark:ch05-径向下紧}
由上述证明，若 $X=\mathbb{R}$，则可取 $K_n=B_n(0)$，即中心在 $0$、半径 $n\in\mathbb{N}$ 的闭球。在此情形，$\varphi$ 可写为
\begin{align*}
\varphi(x)\triangleq \sum_{n=0}^\infty n\,\chi_{(n,n+1]}(|x|)\qquad(x\in X);
\end{align*}
注意它径向、非减且在无穷远发散。
\end{remark}




\end{document}